% ============================================================
%  Activity 8 – Deploy Java App to Cloud (Render)
%  Style mirrors: Docker_Lab_Report.pdf (Activity 6)
% ============================================================
\documentclass[11pt,letterpaper]{article}

% ---------- packages ----------
\usepackage[margin=1in]{geometry}
\usepackage{graphicx}
\usepackage{xcolor}
\usepackage{listings}
\usepackage{booktabs}
\usepackage{tabularx}
\usepackage{hyperref}
\usepackage{fancyhdr}
\usepackage{titlesec}
\usepackage{enumitem}
\usepackage{parskip}
\usepackage{microtype}

% ---------- colours ----------
\definecolor{codebg}{HTML}{F5F5F5}
\definecolor{codeframe}{HTML}{CCCCCC}
\definecolor{keyword}{HTML}{0000FF}
\definecolor{string}{HTML}{A31515}
\definecolor{comment}{HTML}{008000}
\definecolor{gcupurple}{HTML}{522D80}

% ---------- listings ----------
\lstset{
  backgroundcolor=\color{codebg},
  frame=single,
  rulecolor=\color{codeframe},
  basicstyle=\ttfamily\small,
  keywordstyle=\color{keyword}\bfseries,
  stringstyle=\color{string},
  commentstyle=\color{comment}\itshape,
  breaklines=true,
  breakatwhitespace=false,
  showstringspaces=false,
  tabsize=2,
  aboveskip=10pt,
  belowskip=10pt,
  xleftmargin=4pt,
  xrightmargin=4pt,
  numbers=left,
  numberstyle=\tiny\color{gray},
  numbersep=8pt,
}

% ---------- header / footer ----------
\pagestyle{fancy}
\fancyhf{}
\lhead{\textit{Deploy Java App to Cloud (Render)}}
\rhead{\thepage}
\lfoot{Completed: March 31, 2026}
\rfoot{CST-323 Cloud Computing}
\renewcommand{\headrulewidth}{0.4pt}
\renewcommand{\footrulewidth}{0.4pt}

% ---------- section style ----------
\titleformat{\section}{\Large\bfseries}{}{0em}{\thesection\quad}
\titleformat{\subsection}{\large\bfseries}{}{0em}{\thesubsection\quad}

% ---------- hyperlinks ----------
\hypersetup{
  colorlinks=true,
  linkcolor=gcupurple,
  urlcolor=blue!70!black,
  citecolor=gcupurple,
}

% ============================================================
\begin{document}

% ---------- TITLE PAGE ----------
\begin{titlepage}
\centering
\vspace*{2cm}
{\Huge\bfseries Deploy Java App to Cloud Infrastructure\\[6pt]}
{\LARGE Laboratory Completion Report\\[12pt]}
{\Large GCU Cloud-App (Spring Boot + MySQL) on Render\\[24pt]}

\begin{tabular}{@{}rl@{}}
\textbf{Date Completed:}    & March 31, 2026 \\
\textbf{Course:}            & CST-323 Cloud Computing \\
\textbf{Platform:}          & Render.com (Free Web Service) \\
\textbf{Docker Runtime:}    & Eclipse Temurin 17-JRE Alpine \\
\textbf{Spring Boot:}       & v3.2.0 \\
\textbf{MySQL Provider:}    & \textit{<your MySQL host, e.g.\ Railway / Aiven / PlanetScale>} \\
\textbf{Application:}       & cloud-app v1.0.0 \\
\end{tabular}

\vfill
{\small Grand Canyon University --- College of Science, Engineering and Technology}
\end{titlepage}

% ---------- TOC ----------
\tableofcontents
\newpage

% ============================================================
\section{Introduction and Learning Objectives}
% ============================================================

This report documents the deployment of the GCU Cloud-App (Spring Boot + MySQL ``Orders App'') to a cloud hosting platform.
The original assignment targets Oracle Cloud Infrastructure (OCI); however, OCI account creation was blocked during signup (HTTP 403 / ``Forbidden'' on the payment-verification step).
After consulting with the instructor, \textbf{Render.com} was approved as an equivalent deployment target.

Render provides a Platform-as-a-Service (PaaS) model that abstracts away the Virtual Cloud Network (VCN), subnets, and firewall rules that OCI requires.
The core learning outcomes---containerised deployment, cloud database connectivity, environment-variable configuration, and live verification---are preserved.

\subsection{Learning Objectives Achieved}

\begin{itemize}[nosep]
  \item Compare Oracle Cloud's IaaS approach to the PaaS-oriented model used by Render (and AWS, GCP, Azure).
  \item Build and containerise a Spring Boot application using Docker and deploy the image to a cloud service.
  \item Provision and connect to a managed MySQL database from a cloud-hosted container.
  \item Configure environment variables for secure, externalised database credentials.
  \item Diagnose and resolve deployment issues (health-check failures, port binding, datasource overrides).
  \item Evaluate the trade-offs of IaaS-style vs.\ PaaS-style deployment in cloud architecture design.
\end{itemize}

\subsection{System Architecture Overview}

\begin{center}
\begin{tabular}{llll}
\toprule
\textbf{Component} & \textbf{Service} & \textbf{Port} & \textbf{Notes} \\
\midrule
Web Application & Render Web Service (Docker) & 8080 (auto) & Free tier, cold-start enabled \\
Database         & \textit{<MySQL provider>}   & 3306        & Managed MySQL 8              \\
Health Check     & \texttt{/actuator/health}    & ---         & Permitted without login      \\
\bottomrule
\end{tabular}
\end{center}

% ============================================================
\section{IaaS vs.\ PaaS --- Why Render Instead of OCI}
% ============================================================

\subsection{Oracle Cloud's IaaS Approach}

OCI requires developers to manually design the underlying network before deploying any compute or container service:
Virtual Cloud Network (VCN), public and private subnets, Internet Gateway, NAT Gateway, route tables, and Security List ingress/egress rules.
Only after this infrastructure is in place can a Compute VM, Container Instance, or MySQL Database Service be provisioned.

\subsection{Render's PaaS Approach}

Render abstracts the entire networking layer.
Deploying a Docker container requires only:

\begin{enumerate}[nosep]
  \item A GitHub repository with a \texttt{Dockerfile}.
  \item A \textit{Web Service} resource on the Render dashboard.
  \item Environment variables for the datasource.
\end{enumerate}

Render handles TLS termination, DNS, load balancing, and health checking automatically.

\subsection{Side-by-Side Comparison}

\begin{center}
\begin{tabularx}{\textwidth}{lXX}
\toprule
\textbf{Aspect} & \textbf{OCI (IaaS)} & \textbf{Render (PaaS)} \\
\midrule
Network Setup         & Manual VCN, subnets, gateways, route tables & Fully managed; none required \\
Database              & Oracle MySQL Database Service (private subnet) & External MySQL (Railway, Aiven, etc.) \\
Container Hosting     & OCI Container Instances / OCIR                 & Render Web Service (Docker) \\
Firewall Rules        & Manual Security List ingress/egress             & Automatic; outbound open by default \\
TLS / HTTPS           & Manual certificate or Load Balancer             & Automatic via \texttt{*.onrender.com} \\
Cost (free tier)      & Always-Free shapes available                    & Free Web Service (750\,h/month) \\
Deployment Complexity & High --- hands-on DevOps                        & Low --- git-push or image deploy \\
\bottomrule
\end{tabularx}
\end{center}

% ============================================================
\section{OCI Account Issue --- Instructor Approval}
% ============================================================

Multiple attempts to create an Oracle Cloud Free Tier account resulted in an HTTP 403 (``Forbidden'') error during the payment/identity-verification step.
Efforts included:

\begin{itemize}[nosep]
  \item Trying different browsers and clearing cookies.
  \item Using multiple email addresses.
  \item Attempting different payment methods.
  \item Contacting Oracle Support (ticket submitted; 24--48\,h response time).
\end{itemize}

The instructor was notified via email and confirmed that Render is an acceptable substitute.
Screenshots of the Oracle error page and the instructor correspondence are included below.

\textbf{Screenshot 1 --- OCI Signup Error.}
The Oracle Cloud signup page displays a ``Forbidden'' error during identity verification, preventing account creation.

% TODO: replace with your actual screenshot
% \includegraphics[width=\textwidth]{screenshots/oci-forbidden.png}
\begin{center}
\fbox{\parbox{0.8\textwidth}{\centering\vspace{2cm}\textit{Insert screenshot of OCI ``Forbidden'' error here}\vspace{2cm}}}
\end{center}

\textbf{Screenshot 2 --- Instructor Approval.}
Email exchange confirming Render as an approved alternative to OCI for this assignment.

\begin{center}
\fbox{\parbox{0.8\textwidth}{\centering\vspace{2cm}\textit{Insert screenshot of instructor email approval here}\vspace{2cm}}}
\end{center}

% ============================================================
\section{Environment Setup}
% ============================================================

\subsection{Render Account Creation}

A free Render account was created using GitHub OAuth at \url{https://render.com}.
Render was authorised to read the GitHub repository containing the cloud-app source code.

\textbf{Screenshot 3 --- Render Dashboard.}
The Render dashboard after account creation, showing the connected GitHub account.

\begin{center}
\fbox{\parbox{0.8\textwidth}{\centering\vspace{2cm}\textit{Insert screenshot of Render dashboard here}\vspace{2cm}}}
\end{center}

\subsection{GitHub Repository}

The cloud-app source code was pushed to a GitHub repository so Render can build from the \texttt{Dockerfile}.

\begin{lstlisting}[language=bash,caption={Push cloud-app to GitHub}]
cd CST-323-Cloud-Computing/code/topic-2/cloud-app
git init
git remote add origin https://github.com/<username>/<repo>.git
git add .
git commit -m "Initial cloud-app for Render deployment"
git push -u origin main
\end{lstlisting}

% ============================================================
\section{MySQL Database in the Cloud}
% ============================================================

\subsection{Provisioning the Database}

A managed MySQL 8 instance was created on \textit{<your provider>}.
The provider supplies a standard TCP connection string, username, and password.

\textbf{Screenshot 4 --- Cloud MySQL Instance.}
The database provider dashboard showing the MySQL instance is active and the connection details (hostname, port, database name).

\begin{center}
\fbox{\parbox{0.8\textwidth}{\centering\vspace{2cm}\textit{Insert screenshot of MySQL provider dashboard here}\vspace{2cm}}}
\end{center}

\subsection{Schema Initialisation}

Because the application uses Spring Data JDBC (not JPA/Hibernate), the \texttt{USERS} and \texttt{ORDERS} tables must be created manually:

\begin{lstlisting}[language=SQL,caption={Create USERS and ORDERS tables}]
CREATE TABLE IF NOT EXISTS USERS (
  id       INT AUTO_INCREMENT PRIMARY KEY,
  username VARCHAR(100) NOT NULL UNIQUE,
  password VARCHAR(255) NOT NULL,
  role     VARCHAR(50)  NOT NULL,
  enabled  BOOLEAN NOT NULL DEFAULT TRUE
);

CREATE TABLE IF NOT EXISTS ORDERS (
  ID           INT AUTO_INCREMENT PRIMARY KEY,
  ORDER_NUMBER VARCHAR(100),
  PRODUCT_NAME VARCHAR(255),
  PRICE        INT,
  QTY          INT
);
\end{lstlisting}

\textbf{Screenshot 5 --- Tables Created.}
Query result showing both \texttt{USERS} and \texttt{ORDERS} tables exist in the cloud database.

\begin{center}
\fbox{\parbox{0.8\textwidth}{\centering\vspace{2cm}\textit{Insert screenshot of SHOW TABLES result here}\vspace{2cm}}}
\end{center}

% ============================================================
\section{Application Configuration for Render}
% ============================================================

\subsection{Dockerfile}

The multi-stage Dockerfile builds the JAR inside an Alpine JDK container and copies it to a lightweight JRE runtime image.
\texttt{chmod +x ./mvnw} ensures the Maven wrapper is executable on Linux build hosts (Render runs Linux).

\begin{lstlisting}[caption={Dockerfile (project root)}]
FROM eclipse-temurin:17-jdk-alpine AS build
WORKDIR /app
COPY . .
RUN chmod +x ./mvnw && ./mvnw package -DskipTests

FROM eclipse-temurin:17-jre-alpine
WORKDIR /app
COPY --from=build /app/target/*.jar app.jar
EXPOSE 8080
ENTRYPOINT ["java", "-jar", "app.jar"]
\end{lstlisting}

\subsection{application.properties --- Runtime Overrides}

Spring's \texttt{\$\{VAR:default\}} syntax allows every setting to be overridden via environment variables without rebuilding the image:

\begin{lstlisting}[caption={application.properties (key lines)}]
server.port=${PORT:8080}

spring.datasource.url=${SPRING_DATASOURCE_URL:jdbc:mysql://...}
spring.datasource.username=${SPRING_DATASOURCE_USERNAME:admin}
spring.datasource.password=${SPRING_DATASOURCE_PASSWORD:...}
spring.datasource.driver-class-name=com.mysql.cj.jdbc.Driver

spring.jpa.hibernate.ddl-auto=${SPRING_JPA_HIBERNATE_DDL_AUTO:update}

management.endpoints.web.exposure.include=health
management.endpoint.health.probes.enabled=true
\end{lstlisting}

\textbf{Key change:} Actuator health endpoint is now exposed so Render's HTTP health check can reach \texttt{/actuator/health} without authentication.

\subsection{SecurityConfig --- Permit Health Check}

The \texttt{/actuator/health} path was added to Spring Security's permit list so Render's health probe is not redirected to the login page:

\begin{lstlisting}[language=Java,caption={SecurityConfig.java --- health-check permit}]
.authorizeHttpRequests(auth -> auth
    .requestMatchers("/actuator/health", "/actuator/health/**")
        .permitAll()
    .requestMatchers("/users/register", "/users/processRegister",
                     "/users/login").permitAll()
    .requestMatchers("/admin/**").hasRole("ADMIN")
    .anyRequest().authenticated()
)
\end{lstlisting}

% ============================================================
\section{Deploying to Render}
% ============================================================

\subsection{Create a Web Service}

On the Render dashboard: \textbf{New $\rightarrow$ Web Service $\rightarrow$} connect the GitHub repository.

\begin{center}
\begin{tabular}{ll}
\toprule
\textbf{Setting} & \textbf{Value} \\
\midrule
Name             & \texttt{cloud-app} \\
Region           & Oregon (US West) \\
Environment      & Docker \\
Instance Type    & Free \\
Health Check Path & \texttt{/actuator/health} \\
\bottomrule
\end{tabular}
\end{center}

\subsection{Environment Variables}

The following variables were set in the Render dashboard under \textbf{Environment}:

\begin{center}
\begin{tabularx}{\textwidth}{lX}
\toprule
\textbf{Variable} & \textbf{Value (redacted)} \\
\midrule
\texttt{SPRING\_DATASOURCE\_URL}      & \texttt{jdbc:mysql://<host>:3306/<db>?sslMode=REQUIRED\&serverTimezone=UTC} \\
\texttt{SPRING\_DATASOURCE\_USERNAME} & \texttt{<username>} \\
\texttt{SPRING\_DATASOURCE\_PASSWORD} & \texttt{********} \\
\bottomrule
\end{tabularx}
\end{center}

\textit{Note:} Render sets \texttt{PORT} automatically; the application reads it via \texttt{server.port=\$\{PORT:8080\}}.

\textbf{Screenshot 6 --- Render Environment Variables.}
The Render Environment tab showing the \texttt{SPRING\_DATASOURCE\_*} keys are configured (values redacted).

\begin{center}
\fbox{\parbox{0.8\textwidth}{\centering\vspace{2cm}\textit{Insert screenshot of Render Environment tab here}\vspace{2cm}}}
\end{center}

\subsection{Build and Deploy}

Clicking \textbf{Create Web Service} triggers a Docker build on Render's servers.
The build log shows Maven compiling the project, running \texttt{mvnw package}, and creating the runtime image.

\textbf{Screenshot 7 --- Render Build Logs.}
The Render deploy log showing a successful Docker build (Maven compile + JAR packaging) and service startup.

\begin{center}
\fbox{\parbox{0.8\textwidth}{\centering\vspace{2cm}\textit{Insert screenshot of Render build/deploy logs here}\vspace{2cm}}}
\end{center}

% ============================================================
\section{System Verification}
% ============================================================

\subsection{Health Check}

After deployment, the health endpoint confirms the application is running:

\begin{lstlisting}[language=bash,caption={Health check via curl (or browser)}]
$ curl https://<your-app>.onrender.com/actuator/health
{"status":"UP"}
\end{lstlisting}

\textbf{Screenshot 8 --- Health Check.}
Browser or terminal showing \texttt{\{"status":"UP"\}} from the \texttt{/actuator/health} endpoint on the live Render URL.

\begin{center}
\fbox{\parbox{0.8\textwidth}{\centering\vspace{2cm}\textit{Insert screenshot of /actuator/health response here}\vspace{2cm}}}
\end{center}

\subsection{Render Service Status}

\textbf{Screenshot 9 --- Render Service Live.}
The Render dashboard showing the service status as \textbf{Live} with the public URL.

\begin{center}
\fbox{\parbox{0.8\textwidth}{\centering\vspace{2cm}\textit{Insert screenshot of Render service status (Live) here}\vspace{2cm}}}
\end{center}

% ============================================================
\section{CRUD Walkthrough --- Orders4U Web Application}
% ============================================================

With the application live on Render and connected to the cloud MySQL database, the full Create, Read, Update, and Delete lifecycle was exercised through the browser.

\subsection{Authentication --- Login (Spring Security)}

\textbf{Screenshot 10 --- Login Page.}
The login page at \texttt{https://<app>.onrender.com/users/login}.
Spring Security intercepts all unauthenticated requests and redirects to this form.
BCrypt-hashed credentials are verified against the \texttt{USERS} table in the cloud database.

\begin{center}
\fbox{\parbox{0.8\textwidth}{\centering\vspace{2cm}\textit{Insert screenshot of login page here}\vspace{2cm}}}
\end{center}

\subsection{Read --- Home Dashboard (after login)}

\textbf{Screenshot 11 --- Home Dashboard.}
The home page after successful authentication, showing the personalised welcome and navigation to the Orders section.

\begin{center}
\fbox{\parbox{0.8\textwidth}{\centering\vspace{2cm}\textit{Insert screenshot of home dashboard here}\vspace{2cm}}}
\end{center}

\subsection{Read --- Inventory List}

\textbf{Screenshot 12 --- All Orders.}
The inventory page at \texttt{/orders} displaying all current orders retrieved via \texttt{SELECT * FROM ORDERS}.

\begin{center}
\fbox{\parbox{0.8\textwidth}{\centering\vspace{2cm}\textit{Insert screenshot of orders list here}\vspace{2cm}}}
\end{center}

\subsection{Create --- New Order}

\textbf{Screenshot 13 --- New Order.}
The new-order form at \texttt{/orders/newOrder} with sample data entered.
On submit, \texttt{OrdersController} calls \texttt{repository.save()}, issuing an \texttt{INSERT INTO ORDERS}.

\begin{center}
\fbox{\parbox{0.8\textwidth}{\centering\vspace{2cm}\textit{Insert screenshot of new order form here}\vspace{2cm}}}
\end{center}

\subsection{Read --- Single Order Detail}

\textbf{Screenshot 14 --- Order Detail.}
The detail page at \texttt{/orders/showOrder/\{id\}} showing the newly created order.

\begin{center}
\fbox{\parbox{0.8\textwidth}{\centering\vspace{2cm}\textit{Insert screenshot of order detail here}\vspace{2cm}}}
\end{center}

\subsection{Update --- Edit Order}

\textbf{Screenshot 15 --- Edit Order.}
The order after editing (price or quantity changed).
Spring Data JDBC issues \texttt{UPDATE ORDERS SET \ldots\ WHERE ID = ?}.

\begin{center}
\fbox{\parbox{0.8\textwidth}{\centering\vspace{2cm}\textit{Insert screenshot of updated order here}\vspace{2cm}}}
\end{center}

\subsection{Delete --- Remove Order}

\textbf{Screenshot 16 --- Delete Confirmation.}
The delete confirmation dialog followed by the inventory page showing the order has been removed.

\begin{center}
\fbox{\parbox{0.8\textwidth}{\centering\vspace{2cm}\textit{Insert screenshot of delete confirmation here}\vspace{2cm}}}
\end{center}

\subsection{CRUD Summary}

\begin{center}
\begin{tabular}{lllll}
\toprule
\textbf{Operation} & \textbf{Route} & \textbf{HTTP} & \textbf{SQL Issued} & \textbf{Screenshot} \\
\midrule
Login      & \texttt{/users/login}            & POST & SELECT (auth)      & 10 \\
Read All   & \texttt{/orders}                 & GET  & SELECT *           & 12 \\
Create     & \texttt{/orders/newOrder}         & POST & INSERT INTO        & 13 \\
Read One   & \texttt{/orders/showOrder/\{id\}} & GET  & SELECT WHERE ID=?  & 14 \\
Update     & \texttt{/orders/editOrder/\{id\}} & POST & UPDATE WHERE ID=?  & 15 \\
Delete     & \texttt{/orders/deleteOrder/\{id\}} & POST & DELETE WHERE ID=? & 16 \\
\bottomrule
\end{tabular}
\end{center}

% ============================================================
\section{Collaborating with AI}
% ============================================================

% CHOOSE ONE of 1A, 2A, or 3A and fill in your reflection below.
% Delete the other two subsections.

\subsection{3A --- Cross-Cloud Reflection}

\textbf{Prompt:} ``How does deploying a Java application with MySQL on Oracle Cloud compare to similar workflows on AWS (Aurora + ECS), Azure (Database for MySQL + Container Instances), or Google Cloud (SQL + Cloud Run)? Which platform might best suit different project priorities such as cost, performance, or ease of use?''

\textbf{Summary of Findings:}

% TODO: Replace the text below with your own 150--200 word reflection.

\textit{<Your 150--200 word reflection here. Summarise what the AI told you, then add your own analysis. Conclude with which platform you would personally choose for a future project, and why.>}

% ============================================================
\section{Conclusion and Next Steps}
% ============================================================

\subsection{What Was Accomplished}

\begin{enumerate}[nosep]
  \item Documented the OCI account-creation issue and obtained instructor approval for Render as a substitute.
  \item Created a Render Web Service backed by a Docker container built from the cloud-app \texttt{Dockerfile}.
  \item Provisioned a managed MySQL 8 database and created the required \texttt{USERS} and \texttt{ORDERS} tables.
  \item Configured \texttt{SPRING\_DATASOURCE\_*} environment variables on Render for secure credential injection.
  \item Exposed the \texttt{/actuator/health} endpoint and permitted it in Spring Security for Render's health probe.
  \item Verified the live deployment via health check, login, and a full CRUD walkthrough of the Orders application.
\end{enumerate}

\subsection{Suggested Next Steps}

\begin{itemize}[nosep]
  \item \textbf{CI/CD Pipeline:} Add GitHub Actions to build, push to GHCR, and trigger a Render Deploy Hook automatically on each push (Activity~9).
  \item \textbf{Custom Domain:} Configure a custom domain with Render's automatic TLS.
  \item \textbf{Production Database:} Migrate to a production-grade managed MySQL with automated backups.
  \item \textbf{Secrets Management:} Use Render's secret files or a vault service instead of plain environment variables.
  \item \textbf{Monitoring:} Integrate Spring Boot Actuator metrics with an observability platform (Grafana, Datadog).
\end{itemize}

\vfill
\begin{center}
\small CST-323 Cloud Computing \textbullet{} March 31, 2026
\end{center}

\end{document}
